% Language: LaTeX
% =====================================================================================================
%  ejemplos_de_code.tex  --  Una muestra de cada caja de codigo definida en listings.tex
%  Uso (apendices):  \input{secciones/6_apendices/ejemplos_de_code}
%  Requiere: listings.tex (cajas y float 'codigo')
%
%  Entorno          Argumentos             Estilo por defecto
%  Ccode            {barra}                caja_c
%  SVcode           {barra}                caja_sv
%  Bashcode         {barra}                caja_bash
%  Latcode          [estilo]{barra}        caja_c
%  Tabcode          [estilo]{barra}        caja_c
%  Ieeecode         [estilo]{barra}        caja_c
%  TermROGcode      [estilo]{barra}        caja_c_dark
%  TermGGGcode      [estilo]{barra}        caja_c_dark
%  TermLinuxcode    [estilo]{barra}        caja_c_dark
%  MinCcode         [estilo]               caja_c
%  MinCcapcode / MinSVcapcode / MinBashcapcode   (sin argumentos)
%  lstlisting       [style=codigo | codigo_sv, ...]   (formato anterior)
%
%  Patron de uso (el float aporta caption, label e indice):
%    \begin{codigo}
%    \begin{ENTORNO}[estilo]{titulo} ... \end{ENTORNO}
%    \caption{...}
%    \label{lst:...}
%    \end{codigo}
% =====================================================================================================
\section{Ejemplos de código}
\label{sec:ej-codigo}

Este apartado muestra una caja de cada tipo definida en \texttt{listings.tex}. Las cajas solo dibujan el marco; el \texttt{\textbackslash caption}, el \texttt{\textbackslash label} y la numeración «Código X.Y» los aporta el float \texttt{codigo} que las envuelve.

% =====================================================================================================
\subsection{Barra de título: estándar}
\label{subsec:ej-code-estandar}

Barra oscura con un único argumento (el texto de la barra). Un entorno por lenguaje: \texttt{Ccode}, \texttt{SVcode} y \texttt{Bashcode}.

\begin{codigo}
\begin{Ccode}{Ccode: barra oscura, sintaxis C}
#include <stdint.h>

/* Returns the sign bit of an IEEE 754 binary32 word */
static inline uint32_t f32_sign(uint32_t bits) {
    return bits >> 31;   // 1 = negative
}
\end{Ccode}
\caption{Ejemplo con \texttt{Ccode}.}
\label{lst:ej-ccode}
\end{codigo}

\begin{codigo}
\begin{SVcode}{SVcode: barra oscura, sintaxis SystemVerilog}
module counter #(parameter int W = 8) (
  input  logic         clk, rst_n,
  output logic [W-1:0] count
);
  // Synchronous counter with active-low asynchronous reset
  always_ff @(posedge clk or negedge rst_n)
    if (!rst_n) count <= '0;
    else        count <= count + 1;
endmodule
\end{SVcode}
\caption{Ejemplo con \texttt{SVcode}.}
\label{lst:ej-svcode}
\end{codigo}

\begin{codigo}
\begin{Bashcode}{Bashcode: barra oscura, sintaxis bash}
#!/bin/bash
# Build and run a single test
make clean
make all SEED=42
./run_sim.sh | tee sim.log   # keep a copy of the output
\end{Bashcode}
\caption{Ejemplo con \texttt{Bashcode}.}
\label{lst:ej-bashcode}
\end{codigo}

% =====================================================================================================
\subsection{Barra de título: estéticas}
\label{subsec:ej-code-esteticas}

Diseños alternativos con argumentos \texttt{[estilo]\{barra\}}. El estilo de sintaxis es opcional (por defecto C); aquí se usa el estilo por defecto en \texttt{Latcode} y se cambia en \texttt{Tabcode} y \texttt{Ieeecode}.

\begin{codigo}
\begin{Latcode}{Latcode: franja lateral azul}
#include <stdint.h>

/* Returns the sign bit of an IEEE 754 binary32 word */
static inline uint32_t f32_sign(uint32_t bits) {
    return bits >> 31;   // 1 = negative
}
\end{Latcode}
\caption{Ejemplo con \texttt{Latcode} (estilo por defecto \texttt{caja\_c}).}
\label{lst:ej-latcode}
\end{codigo}

\begin{codigo}
\begin{Tabcode}[caja_sv]{Tabcode: pestaña oscura}
module counter #(parameter int W = 8) (
  input  logic         clk, rst_n,
  output logic [W-1:0] count
);
  // Synchronous counter with active-low asynchronous reset
  always_ff @(posedge clk or negedge rst_n)
    if (!rst_n) count <= '0;
    else        count <= count + 1;
endmodule
\end{Tabcode}
\caption{Ejemplo con \texttt{Tabcode} y estilo \texttt{caja\_sv}.}
\label{lst:ej-tabcode}
\end{codigo}

\begin{codigo}
\begin{Ieeecode}[caja_bash]{Ieeecode: barra azul IEEE}
#!/bin/bash
# Build and run a single test
make clean
make all SEED=42
./run_sim.sh | tee sim.log   # keep a copy of the output
\end{Ieeecode}
\caption{Ejemplo con \texttt{Ieeecode} y estilo \texttt{caja\_bash}.}
\label{lst:ej-ieeecode}
\end{codigo}

% =====================================================================================================
\subsection{Terminales oscuros}
\label{subsec:ej-code-terminal}

Comparten los estilos de sintaxis y reglas de fondo oscuro (\texttt{caja\_c\_dark}, \texttt{caja\_sv\_dark} y \texttt{caja\_bash\_dark}) y se diferencian solo en los botones de la barra: círculos de macOS, tres círculos verdes y botón [+] de GNOME.

\begin{codigo}
\begin{TermROGcode}[caja_bash_dark]{TermROGcode: terminal estilo macOS}
cd build
make clean
make all SEED=42   # fixed seed for reproducibility
./run_sim.sh | tee sim.log
\end{TermROGcode}
\caption{Ejemplo con \texttt{TermROGcode} y estilo \texttt{caja\_bash\_dark}.}
\label{lst:ej-termrog}
\end{codigo}

\begin{codigo}
\begin{TermGGGcode}{TermGGGcode: terminal con círculos verdes}
#include <stdint.h>

/* Returns the sign bit of an IEEE 754 binary32 word */
static inline uint32_t f32_sign(uint32_t bits) {
    return bits >> 31;   // 1 = negative
}
\end{TermGGGcode}
\caption{Ejemplo con \texttt{TermGGGcode} (estilo por defecto \texttt{caja\_c\_dark}).}
\label{lst:ej-termggg}
\end{codigo}

\begin{codigo}
\begin{TermLinuxcode}[caja_sv_dark]{TermLinuxcode: terminal estilo GNOME}
module counter #(parameter int W = 8) (
  input  logic         clk, rst_n,
  output logic [W-1:0] count
);
  // Synchronous counter with active-low asynchronous reset
  always_ff @(posedge clk or negedge rst_n)
    if (!rst_n) count <= '0;
    else        count <= count + 1;
endmodule
\end{TermLinuxcode}
\caption{Ejemplo con \texttt{TermLinuxcode} y estilo \texttt{caja\_sv\_dark}.}
\label{lst:ej-termlinux}
\end{codigo}

% =====================================================================================================
\subsection{Sin barra de título}
\label{subsec:ej-code-sinbarra}

Cajas sin barra superior. \texttt{MinCcode} es un marco fino sobre fondo blanco con estilo opcional (por defecto C); \texttt{MinCcapcode}, \texttt{MinSVcapcode} y \texttt{MinBashcapcode} no reciben argumentos y dibujan el marco sobre fondo gris.

\begin{codigo}
\begin{MinCcode}
$#include <stdint.h>

/* Returns the sign bit of an IEEE 754 binary32 word */
static inline uint32_t f32_sign(uint32_t bits) {
    return bits >> 31;   // 1 = negative
}
\end{MinCcode}
\caption{Ejemplo con \texttt{MinCcode} (estilo por defecto \texttt{caja\_c}).}
\label{lst:ej-minc}
\end{codigo}

\begin{codigo}
\begin{MinCcapcode}
#include <stdint.h>
/* Returns the sign bit of an IEEE 754 binary32 word */
static inline uint32_t f32_sign(uint32_t bits) {
    return bits >> 31;   // 1 = negative
}
\end{MinCcapcode}
\caption{Ejemplo con \texttt{MinCcapcode}.}
\label{lst:ej-mincap-c}
\end{codigo}

\begin{codigo}
\begin{MinSVcapcode}
module counter #(parameter int W = 8) (
  input  logic         clk, rst_n,
  output logic [W-1:0] count
);
  // Synchronous counter with active-low asynchronous reset
  always_ff @(posedge clk or negedge rst_n)
    if (!rst_n) count <= '0;
    else        count <= count + 1;
endmodule
\end{MinSVcapcode}
\caption{Ejemplo con \texttt{MinSVcapcode}.}
\label{lst:ej-mincap-sv}
\end{codigo}

\begin{codigo}
\begin{MinBashcapcode}
#!/bin/bash
# Build and run a single test
make clean
make all SEED=42
./run_sim.sh | tee sim.log   # keep a copy of the output
\end{MinBashcapcode}
\caption{Ejemplo con \texttt{MinBashcapcode}.}
\label{lst:ej-mincap-bash}
\end{codigo}

% =====================================================================================================
\subsection{Formato anterior (\texttt{lstlisting})}
\label{subsec:ej-code-anterior}

Formato independiente del float \texttt{codigo}: usa directamente \texttt{lstlisting} con los estilos \texttt{codigo} y \texttt{codigo\_sv}, y la numeración nativa de \texttt{listings}.

\begin{lstlisting}[style=codigo, language=C,
  caption={Ejemplo con el formato anterior (\texttt{style=codigo}).},
  label=lst:ej-legacy-c]
#include <stdint.h>

/* Returns the sign bit of an IEEE 754 binary32 word */
static inline uint32_t f32_sign(uint32_t bits) {
    return bits >> 31;   // 1 = negative
}
\end{lstlisting}

\begin{lstlisting}[style=codigo_sv,
  caption={Ejemplo con el formato anterior (\texttt{style=codigo\_sv}).},
  label=lst:ej-legacy-sv]
module counter #(parameter int W = 8) (
  input  logic         clk, rst_n,
  output logic [W-1:0] count
);
  // Synchronous counter with active-low asynchronous reset
  always_ff @(posedge clk or negedge rst_n)
    if (!rst_n) count <= '0;
    else        count <= count + 1;
endmodule
\end{lstlisting}